% !TEX root = Projektdokumentation.tex
\section{Anhang}

\subsection{Detaillierte Zeitplanung}
\label{app:Zeitplanung}
\tabelleAnhang{ZeitplanungKomplett}

\subsection{Lastenheft (Auszug)}
\label{app:Lastenheft}

% TODO: Auszug aus dem Lastenheft / den Anforderungen ergänzen,
% z.\,B. die Punkte aus dem Projektantrag (2.2 Soll-Konzept,
% 2.3 zu treffende Entscheidungen, 2.5 Einschränkungen).


\subsection{Konfigurations-Listings}
\label{app:Listings}
Die folgenden Auszüge zeigen die produktiven Konfigurations\-dateien
auf \texttt{odcf-fs04} (Stand der Blueprint-Instanz). Vertrauliche
Bestandteile (DKFZ-LDAP-Bind, Distinguished Names, Filter,
LDAP-Server) sind durch \texttt{XXX} ersetzt.

\begin{lstlisting}[caption={\texttt{/etc/network/interfaces} -- LACP-Bond aus 2$\times$10\,GbE},label=lst:interfaces]
auto lo
iface lo inet loopback

iface nic0 inet manual
iface nic1 inet manual
iface nic2 inet manual
iface nic3 inet manual

auto nic4
iface nic4 inet manual

auto nic5
iface nic5 inet manual

auto bond0
iface bond0 inet static
        address 10.131.196.78/22
        gateway 10.131.196.97
        bond-slaves nic4 nic5
        bond-miimon 100
        bond-mode 802.3ad
        bond-xmit-hash-policy layer3+4
        bond-lacp-rate fast

source /etc/network/interfaces.d/*
\end{lstlisting}

\begin{lstlisting}[caption={\texttt{/etc/proxmox-backup/datastore.cfg}},label=lst:datastore]
datastore: backup
        path /mnt/datastore/backup
\end{lstlisting}

\begin{lstlisting}[caption={\texttt{/etc/proxmox-backup/domains.cfg} (DKFZ-spezifische Daten geschwärzt)},label=lst:ldap]
pam: pam
        comment Linux PAM standard authentication
        type pam

pbs: pbs
        comment Proxmox Backup authentication server
        type pbs

ad: dkfz-heidelberg.de
        base-dn XXX
        bind-dn XXX
        default true
        filter XXX
        mode ldaps
        port 636
        server1 XXX
        sync-attributes email=mail,firstname=givenName,lastname=sn
        user-classes user
\end{lstlisting}


\subsection{Betriebsdokumentation (Auszug)}
\label{app:BetriebsDoku}

% TODO: Auszug aus der Betriebsdokumentation für das ODCF-Team einfügen,
% z.\,B.:
%   - Installation und Grundkonfiguration des PBS
%   - Datastore-Konfiguration (ZFS-Pool, Quotas)
%   - Anbindung an den PVE-Cluster (API-Token, Fingerprint)
%   - Beispiel-Backup-Job mit Schedule und Retention
%   - Restore-Verfahren (einzelne Datei, ganze VM)

\clearpage

\clearpage
\subsection{Gesamtarchitektur}
\label{app:Architektur}
Die folgende Übersicht zeigt die Hauptkomponenten der Lösung: den
\ac{PVE}-Cluster als Backup-Quelle, den neuen \ac{PBS}
\texttt{odcf-fs04} mit eigenem Datastore, die flankierenden Dienste
(\ac{DKFZ}-\ac{LDAP}, \ac{Prometheus}, Bitwarden) sowie die parallel
weiterlaufende \ac{TSM}-/Tape-Sicherung.

\begin{figure}[htbp]
    \centering
    \includegraphics[width=0.95\linewidth,keepaspectratio]{architecture}
    \caption{Komponenten- und Datenflüsse der PBS-Blueprint-Lösung}
    \label{fig:Architektur}
\end{figure}
\clearpage

\subsection{Ablauf eines Backup-Jobs}
\label{app:BackupSeq}
Der folgende Sequenzablauf veranschaulicht, wie ein geplanter
Backup-Job vom \ac{PVE}-Hypervisor ausgelöst, am \ac{PBS}
authentifiziert und im Datastore inkrementell abgelegt wird.

\begin{figure}[htbp]
    \centering
    \includegraphics[width=0.85\linewidth,keepaspectratio]{backup-sequence}
    \caption{Ablauf eines Backup-Jobs (PVE → PBS) mit anschließendem Verify}
    \label{fig:BackupSeq}
\end{figure}
\clearpage

\subsection{Ablauf eines File-Restores}
\label{app:RestoreSeq}
Beim File-Restore startet der \ac{PBS} eine temporäre Restore-\ac{VM},
mountet den ausgewählten Snapshot read-only und stellt die Datei
über die Web-Oberfläche zur Verfügung.

\begin{figure}[htbp]
    \centering
    \includegraphics[width=0.85\linewidth,keepaspectratio]{restore-sequence}
    \caption{File-Restore aus dem PBS über die PVE-Weboberfläche}
    \label{fig:RestoreSeq}
\end{figure}
\clearpage

\subsection{PVE-Cluster-Topologie}
\label{app:ClusterTopologie}
Das folgende Diagramm zeigt die Netzwerk-Topologie des
\ac{PVE}-Clusters der \ac{ODCF} mit den sechs Hypervisor-Knoten
(verteilt auf 2.\,OG, 4.\,OG und Rechenzentrum), den jeweiligen
Switch-Paaren je Etage, den \ac{LACP}-Bonds für das Ceph- und
\ac{PVE}-/Management-Netz sowie der Anbindung des neuen
\ac{PBS} \texttt{odcf-fs04}.
Quelle: Mermaid-Diagramm \texttt{mermaid/pve-cluster.mmd} im Repo.

\begin{landscape}
\begin{figure}[htbp]
    \centering
    \includegraphics[width=\linewidth,height=0.85\textheight,keepaspectratio]{pve-cluster}
    \caption{PVE-Cluster-Topologie der ODCF (Ceph: durchgezogen, PVE/MGMT: gestrichelt)}
    \label{fig:ClusterTopologie}
\end{figure}
\end{landscape}
\clearpage

% Cluster-Topologie liegt als Mermaid-Source in mermaid/pve-cluster.mmd
% und gerendert als Bilder/pve-cluster.{png,pdf}.
